\section{Background}
\label{sec:background}

\subsection{Multi-scale MCMC (G.11)}

Multi-scale MCMC was introduced in G.11~\citep{dellaLibera2026multiscale} as
a strategy to improve mixing of redistricting Markov chains.
Standard \recom{}~\citep{deford2021} proposes changes by merging two adjacent
districts and re-splitting via a spanning tree.
The proposals are local: two districts move simultaneously, but the change in
the overall plan is small.
For states with many districts and complex county geography, this leads to
slow mixing: the chain requires many steps to explore configurations that
differ substantially from the starting plan.

\ms{} addresses this by introducing a \emph{coarse level}.
Tracts are aggregated into their containing counties; the county-level
adjacency graph is used to build a coarser redistricting problem.
A coarse move proposes a redistricting at the county level, then projects
back to the tract level using a boundary-swap rebalancer with a budget of
200 iterations.
If the rebalancer succeeds in restoring population balance within tolerance,
the projected plan is accepted; otherwise, the coarse move is rejected and
the chain remains at its current plan.

\paragraph{Seeding.}
\ms{} uses the \texttt{MSC\_STEP\_} prefix for deterministic step seeds:
\[
  \text{step\_seed}(\text{base\_seed}, \text{step}, \text{chain\_idx})
  = \text{SHA-256}(
    \texttt{"MSC\_STEP\_"} \;\|\; \text{step}_{64\text{le}}
    \;\|\; \texttt{"\_"} \;\|\; \text{chain\_idx}_{32\text{le}}
    \;\|\; \texttt{"\_"} \;\|\; \text{base\_seed}_{64\text{le}}
  )_{[0..8]\text{ le}}.
\]
This formula is inherited unchanged by \ams{}.

\paragraph{Coarse balance tolerance.}
\ms{} uses a relaxed population tolerance for the coarse level,
$\delta_\text{coarse} = 2\text{--}3 \times \delta_\text{fine}$, to allow the
rebalancer more latitude when projecting county-level plans to the tract level.
\ams{} uses $\delta_\text{coarse} = 3 \times \delta_\text{fine}$ (the more
relaxed default) to avoid over-rejection during early adaptation when $\alpha$
may be miscalibrated.

\subsection{Robbins-Monro Stochastic Approximation}

\citet{robbins1951} introduced the following problem.
Let $h: \mathbb{R} \to \mathbb{R}$ be an unknown function with a unique root
$\alpha^*$ (i.e., $h(\alpha^*) = 0$).
At each round $t$, we can observe $h(\alpha_t) + \varepsilon_t$ where
$\varepsilon_t$ is mean-zero noise.
The Robbins-Monro procedure updates:
\[
  \alpha_{t+1} = \alpha_t - \gamma_t \bigl(h(\alpha_t) + \varepsilon_t\bigr),
\]
with step sizes $\{\gamma_t\}$ satisfying:
\begin{equation}
  \sum_{t=1}^\infty \gamma_t = \infty
  \qquad\text{and}\qquad
  \sum_{t=1}^\infty \gamma_t^2 < \infty.
  \label{eq:rm-conditions}
\end{equation}

\begin{theorem}[Robbins-Monro, 1951]
  \label{thm:rm}
  Under~\eqref{eq:rm-conditions} and the conditions that $h$ is monotone
  increasing, $\mathrm{E}[\varepsilon_t^2]$ is bounded, and the iterates
  $\{\alpha_t\}$ remain bounded, $\alpha_t \to \alpha^*$ in mean square.
\end{theorem}

In our application, $h(\alpha) = \mathrm{E}[\text{coarse accept rate}(\alpha)] - \tau$
where $\tau = 0.30$ is the target acceptance rate.
The ``noise'' $\varepsilon_t$ is the difference between the observed acceptance
rate in the current window and its expectation.
The monotonicity condition --- that the expected coarse acceptance rate is
increasing in $\alpha$ --- is empirically satisfied for typical redistricting
county graphs but is not formally proved.
We discuss this caveat in Section~\ref{sec:convergence}.

\paragraph{Step-size schedule.}
The schedule $\gamma_t = \gamma_0 / \sqrt{t}$ satisfies~\eqref{eq:rm-conditions}:
\begin{align*}
  \sum_{t=1}^\infty \frac{\gamma_0}{\sqrt{t}} &= \infty
  \quad \text{(divergent $p$-series, $p = 1/2 \le 1$)}, \\
  \sum_{t=1}^\infty \frac{\gamma_0^2}{t} &= \infty
  \quad \text{(harmonic series — borderline; but } \gamma_t^2 = \gamma_0^2/t
  \text{ and } \sum 1/t \text{ diverges)}.
\end{align*}
We note that $\gamma_t = \gamma_0/t$ would give $\sum \gamma_t^2 = \gamma_0^2 \pi^2/6 < \infty$
and $\sum \gamma_t = \gamma_0 \cdot \infty$, which also satisfies the conditions and converges
faster in the early iterations.
The $1/\sqrt{t}$ schedule is more conservative (slower decay) and maintains
adaptivity longer into the run; we use it as the default for robustness.

\subsection{Adaptive MCMC Theory}

Adaptive MCMC --- modifying the proposal kernel online during a run ---
has been studied since the early 2000s.
\citet{haario2001} introduced the Adaptive Metropolis (AM) algorithm, which
updates the covariance of a Gaussian proposal using the empirical covariance
of accepted samples.
\citet{andrieu2008} provide a tutorial covering convergence conditions and
implementation.

The standard theoretical concern with adaptive MCMC is the \emph{diminishing
adaptation} condition: the amount of adaptation must decrease to zero over
time, to ensure the chain asymptotically satisfies detailed balance with
respect to the target distribution.
The Robbins-Monro step-size schedule $\gamma_t \to 0$ ensures diminishing
adaptation: after many adaptation rounds, $\alpha$ changes by only a tiny
amount per round, and the chain is effectively non-adaptive in the long run.

A second concern is \emph{containment}: the adapted chain must not get stuck
in a region where adaptation drives it away from the target.
The clip $\alpha \in [0.05, 0.95]$ ensures containment by keeping the chain
in a regime where both fine and coarse moves occur with positive probability.

Unlike AM-style adapters that target a specific acceptance probability for a
single proposal kernel, \ams{} targets an acceptance rate for the
\emph{coarse kernel specifically}, while the fine kernel runs unmodified.
This makes the adaptation simpler: there is a single scalar parameter $\alpha$
to tune, the monotonicity assumption is straightforward, and convergence is
immediate from the classical Robbins-Monro result.

\subsection{The Three-Layer Compositor}

The redistricting platform uses a three-layer compositor
(Structure / WeightsOverride / SeedCompositor).
\ams{} is a \emph{SeedCompositor} strategy (Search layer), replacing
\texttt{--search multiscale} with \texttt{--search multiscale-adaptive}.
It is orthogonal to the Structure layer (\geosec{}, \areasec{}, prime-factor)
and the WeightsOverride layer (geographic, county, VRA-aligned).
